% Part: first-order-logic
% Chapter: lindstrom
% Section: introduction

\documentclass[../../../include/open-logic-section]{subfiles}

\begin{document}

\olfileid{mod}{lin}{int}
\section{Pambuka}

Ing bab iki kita arep mbuktekake karakterisasi Lindstr\"om
tumrap logika orde kapisan: kanthi sawatara syarat tambahan,
iki logika maksimal kang Teorema Kekompakan lan Teorema
L\"owenheim--Skolem Mudhun tetep lumaku
(\olref[fol][com][com]{thm:compactness} lan
\olref[fol][com][dls]{thm:downward-ls}). Kaping pisan, kita
perlu menehi karakterisasi kang luwih umum tumrap kelas
logika kang dadi cakupan teorema iki. Kita bakal mbatesi
panliten iki marang basa \emph{relasional}, yaiku basa kang
mung ngemot !!{predicate}s lan konstanta individu, tanpa
!!{function}s.

\end{document}
